\section{A perturbation with uniform correlation}
\label{sec:perturbation}

We seek perturbations that tend to zero in $2$-norm while producing
a fixed first-order correlation along words of increasing length.
The nonlinear error must remain bounded independently of that length.
For this purpose we first record the length-one case of Haagerup's
inequality \cite[Lemma~1.4]{Haagerup1979}, together with its elementary
proof.

\begin{lemma}\label{lem:length-one-convolution}
Let $F$ be a free group with a specified free basis, let $\mathcal L$
be the set of basis elements and their inverses, and let $\lambda$
be the left regular representation on $\ell^2(F)$.
For every finitely supported family $(b_\gamma)_{\gamma\in\mathcal L}$,
\[
  \norm{\sum_{\gamma\in\mathcal L}b_\gamma\lambda(\gamma)}
  \le 2\left(\sum_{\gamma\in\mathcal L}|b_\gamma|^2\right)^{1/2}.
\]
\end{lemma}

\begin{proof}
For a letter $a$, let $B_a$ be the orthogonal projection onto the
closed span of reduced words beginning with $a$.  Split
\[
  \lambda(\gamma)
  =\lambda(\gamma)(1-B_{\gamma^{-1}})
   +\lambda(\gamma)B_{\gamma^{-1}}.
\]
The first summand has range in the range of $B_\gamma$.
These ranges are mutually orthogonal, so for $\eta\in\ell^2(F)$,
\begin{align*}
  \norm{\sum_\gamma b_\gamma\lambda(\gamma)
                (1-B_{\gamma^{-1}})\eta}^2
  &=\sum_\gamma |b_\gamma|^2
                 \norm{(1-B_{\gamma^{-1}})\eta}^2\\
  &\le\left(\sum_\gamma |b_\gamma|^2\right)\norm{\eta}^2.
\end{align*}
For the second summands, the triangle inequality and
Cauchy--Schwarz give
\begin{align*}
  \norm{\sum_\gamma b_\gamma\lambda(\gamma)
                    B_{\gamma^{-1}}\eta}
  &\le\sum_\gamma |b_\gamma|\norm{B_{\gamma^{-1}}\eta}\\
  &\le\left(\sum_\gamma|b_\gamma|^2\right)^{1/2}\norm{\eta},
\end{align*}
because the projections $B_{\gamma^{-1}}$ are mutually orthogonal.
Adding the two estimates proves the lemma.
\end{proof}

\begin{proposition}\label{prop:perturbation}
Let $P\subset M$ be an irreducible inclusion of $\mathrm{II}_1$ factors.
Given $u,z\in\mathcal U(M)$, $\xi\in M$, and $0<\alpha<1/2$,
for every integer $\ell\ge1$ there are
$\widetilde u\in\mathcal U(M)$ and
$\widetilde t_0,\ldots,\widetilde t_\ell\in\mathcal U(P)$ such that,
with
\[
  \widetilde W=\widetilde t_0
     \prod_{k=1}^{\ell}(\widetilde u\,\widetilde t_k),
\]
one has
\begin{equation}\label{eq:perturbation-bounds}
\begin{aligned}
  \norm{\widetilde u-u}_2
    &\le \frac{\alpha}{\sqrt{2\ell}}+\frac1\ell,\\
  \left|\tau(\xi^*\widetilde W)
           -\frac\alpha2\tau(\xi^*z)\right|
    &\le \norm{\xi}_2\frac{4\alpha^2}{1-2\alpha}+\frac1\ell.
\end{aligned}
\end{equation}
All products indexed by $k$ are taken in increasing order from left
to right.  The unitaries may depend on $\ell$.
\end{proposition}

\begin{proof}
Fix $\ell$ and work first in $\mathbf M=M^\omega$, with
$\mathbf Q=P^\omega$.  We use $\tau$ also for the ultrapower trace.

\smallskip\noindent
\emph{Free variables and prefixes.}
Set $D_0=C^*(1,u,z,\xi)$.  By
Lemma~\ref{lem:free-haar-extension}, choose successively
$t_j\in\mathcal U(\mathbf Q)$ Haar and free from $D_j$, and put
$D_{j+1}=C^*(D_j,t_j)$, for $0\le j\le\ell$.
Each of these $C^*$-algebras is norm separable.
Define
\[
  s_1=t_0,\qquad s_{j+1}=u t_j\quad(1\le j\le\ell),\qquad
  g_k=s_1\cdots s_k\quad(1\le k\le\ell),
\]
and
\[
  W=g_\ell s_{\ell+1}
    =t_0\prod_{k=1}^{\ell}(u t_k).
\]
Thus $g_k$ is the prefix immediately preceding the $k$th occurrence
of $u$ in $W$.  Since
$D_j=C^*(D_0,s_1,\ldots,s_j)$, Lemma~\ref{lem:haar-absorption}
shows that $s_{j+1}$ is Haar and free from $D_j$.
In particular, $s_1,\ldots,s_\ell$ are free Haar unitaries.
Every nonidentity reduced word in them has trace zero, so the group
they generate can be identified with the free group $G$ on these
generators.
Since $W=(g_\ell u)t_\ell$ and $g_\ell u\in\mathcal U(D_\ell)$,
the same lemma shows that $W$ is Haar and free from $D_\ell$.
Applying it to $W^*$ and $z\in\mathcal U(D_\ell)$ also shows that
\[
  y=zW^*
\]
is Haar and free from $D_\ell$.

For $p\in G$ set $Y_p=pyp^{-1}$.  These form a free family of
Haar unitaries.  Indeed, if $a_1,\ldots,a_r$ are nonzero integers
and successive labels $p_1,\ldots,p_r$ are distinct, then
\[
  Y_{p_1}^{a_1}\cdots Y_{p_r}^{a_r}
   =p_1y^{a_1}(p_1^{-1}p_2)y^{a_2}\cdots
      (p_{r-1}^{-1}p_r)y^{a_r}p_r^{-1}.
\]
Each internal ratio is a nonidentity element of $G$ and hence
is centered in $D_\ell$.  Equation~\eqref{eq:free-endpoints}
gives zero trace.  It follows, by applying this to the reduction
of $v^*w$, that distinct reduced words in the $Y_p$ are orthonormal
in $L^2(\mathbf M,\tau)$, including the empty word.

A left perturbation of $u$ at its $k$th occurrence in $W$ is
conjugated by $g_k$ when moved to the front of the word.  We therefore
use the conjugates $g_j^{-1}yg_j$: conjugating the $j$th one by
$g_j$ returns $y$.  Taking imaginary parts will give a self-adjoint
perturbation, while freeness makes their average small in $2$-norm.
Put $\operatorname{Im}(b)=(b-b^*)/(2i)$ and define
\[
  h=\frac\alpha\ell\sum_{j=1}^{\ell}
                \operatorname{Im}(g_j^{-1}yg_j),\qquad
  u'=e^{ih}u,\qquad
  W'=t_0\prod_{k=1}^{\ell}(u't_k).
\]
The $g_j^{-1}$ are distinct in $G$, so the summands defining $h$
are orthogonal centered imaginary parts of free Haar unitaries.
Each has squared $2$-norm $1/2$.  Consequently,
\[
  h=h^*,\qquad \norm{h}\le\alpha,\qquad
  \norm{h}_2=\frac\alpha{\sqrt{2\ell}}.
\]
Write $H_k=g_khg_k^{-1}$, $p_{kj}=g_kg_j^{-1}$,
and $c=\alpha/(2\ell)$.  Cancelling consecutive prefixes gives
\begin{equation}\label{eq:perturbation-prefix}
  W'=\left(\prod_{k=1}^{\ell}e^{iH_k}\right)W,
  \qquad
  iH_k=c\sum_{j=1}^{\ell}
                  (Y_{p_{kj}}-Y_{p_{kj}}^{-1}).
\end{equation}
For the first identity, use
$e^{iH_k}=g_ke^{ih}g_k^{-1}$ and
$g_k^{-1}g_{k+1}=ut_k$ for $1\le k<\ell$; the product order is
therefore essential.

Let $S=\{p_{kj}:1\le k,j\le\ell\}$ and let $n_p$ count the
ordered pairs $(k,j)$ with $p_{kj}=p$.
For the identity $e$ of $G$, one has $n_e=\ell$.
If $k\ne j$, the word
\[
  p_{kj}=s_1\cdots s_k s_j^{-1}\cdots s_1^{-1}
\]
is reduced, since its joining letters have different indices.
Its initial positive segment and its final negative segment
determine $k$ and $j$.  All off-diagonal pairs therefore give
distinct nonidentity words, and
\begin{equation}\label{eq:perturbation-multiplicity}
  \sum_{p\in S}n_p^2=\ell^2+\ell(\ell-1)=2\ell^2-\ell.
\end{equation}

\smallskip\noindent
\emph{The first-order correlation.}
Put $C_1=\sum_{k=1}^{\ell}iH_k$.  Each diagonal term
$cY_{p_{kk}}W$ equals $cz$, because $p_{kk}=e$ and $Y_eW=z$.
We show that all other linear terms vanish when paired with $\xi$.
Equation~\eqref{eq:free-endpoints} for $W$ and $D_\ell$ gives
$\tau(\xi^*W)=0$.  For $p\ne e$, the element $p^{-1}$ is
centered in $D_\ell$, so
\[
  \tau(\xi^*Y_pW)
   =\tau(\xi^*p z W^*p^{-1}W)=0.
\]
For every $p\in G$, one also has
\[
  \tau(\xi^*Y_p^{-1}W)
   =\tau(\xi^*p W z^*p^{-1}W)=0.
\]
Indeed, split the internal factor $z^*p^{-1}\in D_\ell$ into
its centered and scalar parts.  The first contribution vanishes
by \eqref{eq:free-endpoints}; the second is a scalar multiple of
$\tau(\xi^*pW^2)$ and vanishes by the same formula.
It follows that
\[
  \tau\bigl(\xi^*(1+C_1)W\bigr)
   =c n_e\tau(\xi^*z)=\frac\alpha2\tau(\xi^*z).
\]

\smallskip\noindent
\emph{The uniform remainder estimate.}
It remains to bound the contribution of
$\prod_k e^{iH_k}-1-C_1$ uniformly in $\ell$.
Group the ordered product of exponentials by its original degree:
\[
  C_d=\sum_{\substack{m_1+\cdots+m_\ell=d\\m_k\ge0}}
       \frac{(iH_1)^{m_1}\cdots(iH_\ell)^{m_\ell}}
            {m_1!\cdots m_\ell!}.
\]
The term $C_0$ is $1$, and the degree-one term is the $C_1$ above.
For fixed $\ell$, the expansion is absolutely convergent in
operator norm, because
\[
  \sum_{d\ge0}\norm{C_d}
  \le\prod_{k=1}^{\ell}
          \sum_{m\ge0}\frac{\norm{H_k}^m}{m!}
  =\exp\left(\sum_{k=1}^{\ell}\norm{H_k}\right)
  \le e^{\ell\alpha}.
\]
The Cauchy product may thus be grouped in this way, and its sum is
$\prod_k e^{iH_k}$.
We will prove the bound
\begin{equation}\label{eq:perturbation-remainder}
  \norm{\prod_{k=1}^{\ell}e^{iH_k}-1-C_1}_2
  \le\frac{4\alpha^2}{1-2\alpha},
\end{equation}
which, unlike the preceding norm-convergence bound, is uniform
in $\ell$.

Let $F_S$ be the formal free group with generators $T_p$, $p\in S$.
In particular $T_e$ is a generator, not the identity of $F_S$.
Let $\Phi:\mathbb C[F_S]\to\mathbf M$ be the homomorphism
$\Phi(T_p)=Y_p$.
By the orthonormality proved above, for every finite group-ring sum
$f=\sum_{w\in F_S}f(w)w$ one has
\begin{equation}\label{eq:perturbation-coefficient-isometry}
  \norm{\Phi(f)}_2^2=\sum_{w\in F_S}|f(w)|^2.
\end{equation}
Define
\[
  B_k=c\sum_{j=1}^{\ell}
                (T_{p_{kj}}-T_{p_{kj}}^{-1}),\qquad
  A_k=c\sum_{j=1}^{\ell}
                (T_{p_{kj}}+T_{p_{kj}}^{-1}),
\]
and
\[
  A=\sum_{k=1}^{\ell}A_k
    =c\sum_{p\in S}n_p(T_p+T_p^{-1}).
\]
For a group-ring element $f$, write $|f|$ for its coefficientwise
absolute value, and compare nonnegative coefficient functions
coefficientwise.  Group multiplication gives
\[
  |fg|(w)=\left|\sum_{ab=w}f(a)g(b)\right|
      \le\sum_{ab=w}|f(a)||g(b)|=(|f|\,|g|)(w).
\]
This comparison already includes all free reductions and all
contributions collected at the same resulting word.
Thus the formal element
\[
  \mathcal C_d=
     \sum_{\substack{m_1+\cdots+m_\ell=d\\m_k\ge0}}
        \frac{B_1^{m_1}\cdots B_\ell^{m_\ell}}
             {m_1!\cdots m_\ell!},
  \qquad \Phi(\mathcal C_d)=C_d,
\]
satisfies
\begin{equation}\label{eq:perturbation-majorization}
  |\mathcal C_d|
   \le\sum_{m_1+\cdots+m_\ell=d}
            \frac{A_1^{m_1}\cdots A_\ell^{m_\ell}}
                 {m_1!\cdots m_\ell!}
   \le A^d.
\end{equation}
For the last inequality, each tuple $(m_1,\ldots,m_\ell)$
corresponds to one nondecreasing sequence of row indices in the
distributive expansion of $(A_1+\cdots+A_\ell)^d$.
Its factorial weight is at most one, and every other summand has
nonnegative coefficients.  No commutation of the $A_k$ is used.
The index $d$ records the degree before reduction; resulting words
may have shorter length, including zero.

Let $\lambda$ now denote the left regular representation of $F_S$,
and let $\delta_{\mathrm{id}}$ be the unit vector at its identity.
By \eqref{eq:perturbation-coefficient-isometry} and
\eqref{eq:perturbation-majorization},
\[
  \norm{C_d}_2
   \le\norm{A^d}_{\ell^2(F_S)}
   =\norm{\lambda(A)^d\delta_{\mathrm{id}}}_{\ell^2(F_S)}
   \le\norm{\lambda(A)}^d.
\]
This uses only the coefficient $2$-norm identity
\eqref{eq:perturbation-coefficient-isometry}, with no comparison
of operator norms between the two representations.
Lemma~\ref{lem:length-one-convolution} and
\eqref{eq:perturbation-multiplicity} give
\[
  \norm{\lambda(A)}
    \le 2\left(2c^2\sum_{p\in S}n_p^2\right)^{1/2}
    =2\alpha\left(1-\frac1{2\ell}\right)^{1/2}
    \le2\alpha<1.
\]
Consequently $\norm{C_d}_2\le(2\alpha)^d$.
The operator-norm convergence established above implies convergence
of the same partial sums in $2$-norm.  Taking their tails and using
the triangle inequality therefore yields
\[
  \norm{\prod_{k=1}^{\ell}e^{iH_k}-1-C_1}_2
  \le\sum_{d\ge2}(2\alpha)^d
  =\frac{4\alpha^2}{1-2\alpha}.
\]
In particular, orthogonality between different degrees is not
required.  This proves \eqref{eq:perturbation-remainder}.

Combining the first-order correlation with
\eqref{eq:perturbation-prefix}, \eqref{eq:perturbation-remainder},
and Cauchy--Schwarz gives
\begin{equation}\label{eq:perturbation-ultrapower-correlation}
  \left|\tau(\xi^*W')-\frac\alpha2\tau(\xi^*z)\right|
  \le\norm{\xi}_2\frac{4\alpha^2}{1-2\alpha}.
\end{equation}

\smallskip\noindent
\emph{Lifting for the fixed word length.}
Each $t_j$ has a representative consisting of unitaries in $P$.
To see this, start with a bounded representative $(b_n)$ of any
unitary in $P^\omega$.  The polar part of $b_n$ extends to a
unitary $w_n\in P$, since its initial and final complementary
projections have equal trace in the finite factor $P$.
Then $b_n=w_n|b_n|$ and functional calculus gives
\[
  \norm{b_n-w_n}_2=\norm{|b_n|-1}_2
     \le\norm{b_n^*b_n-1}_2\longrightarrow_\omega0.
\]
Here the inequality follows from
$|\sqrt t-1|\le|t-1|$ for $t\ge0$.

Choose unitary representatives $t_{j,n}\in P$ for the finitely
many $t_j$.  Using the constant representatives of $u,z$, form
$g_{k,n},W_n,y_n,h_n,u_n',W_n'$ by exactly the formulas above.
For this fixed $\ell$ these are bounded sequences, with
$h_n=h_n^*$ and $\norm{h_n}\le\alpha$.
Bounded algebraic operations and continuous functional calculus
in the quotient show that they represent the corresponding
ultrapower elements.  In particular,
\[
  \lim_{n\to\omega}\norm{h_n}_2
       =\frac\alpha{\sqrt{2\ell}},\qquad
  \lim_{n\to\omega}\tau(\xi^*W_n')=\tau(\xi^*W').
\]
Moreover $\norm{u_n'-u}_2\le\norm{h_n}_2$, by the scalar
inequality $|e^{it}-1|\le|t|$ and functional calculus.
Choose one index $n$ for which
\[
  \norm{h_n}_2<\frac\alpha{\sqrt{2\ell}}+\frac1\ell,
  \qquad
  |\tau(\xi^*W_n')-\tau(\xi^*W')|<\frac1\ell.
\]
Both conditions hold on sets belonging to $\omega$, so their
intersection is nonempty.  Set $\widetilde u=u_n'$ and
$\widetilde t_j=t_{j,n}$.  Equation
\eqref{eq:perturbation-ultrapower-correlation} proves
\eqref{eq:perturbation-bounds}.  This choice is made separately
for each $\ell$, as the proposition permits.
\end{proof}
